\section{Introduction}

Indoor mobile robots often operate in environments where individual sensing modalities are intermittently informative. Range-bearing or scan-based perception can be degraded by weak geometric structure, while radio or beacon measurements can be affected by range limits, wall attenuation, and changing visibility. The repository studied here isolates this sensor-fusion problem in a reproducible synthetic setting: a ground robot moves through known 2D wall-segment maps, receives noisy odometry-like prediction inputs, and fuses intermittent LiDAR-derived pose measurements with range-only beacon measurements.

The implemented problem is localization rather than SLAM. The map and beacon locations are known, and the estimator does not estimate wall geometry or beacon positions. The main question is whether a beacon range covariance adapted from robot-side measurements, specifically measured range and RSSI, improves an error-state Kalman filter relative to a fixed beacon covariance. A separate LiDAR-only and LiDAR-plus-beacon comparison quantifies the role of scan-derived pose corrections.

The state representation follows the manifold-aware structure used in modern robotics estimation libraries such as GTSAM \citep{dellaert2022gtsam,gtsamConcepts,dellaert2012gtsam,dellaertKaess2017factor}. Although the code does not call GTSAM, it uses the same conceptual separation between a nominal state on a nonlinear space and a local vector perturbation used for linearization. This is relevant because robot orientation belongs to a rotation group, and applying heading corrections by unconstrained vector addition can obscure the distinction between global state variables and local tangent coordinates. The filter therefore stores the nominal estimate as a matrix NavState and applies prediction and correction increments through a NavState exponential map.
